\documentclass[a4paper, 10pt]{article}
\usepackage{scrextend}
\changefontsizes{9pt}
% Per-subject config for this repo. See vendor/template/course-config.tex.example
% for what each macro is used for.

\newcommand{\CourseCode}{2110506}
\newcommand{\CourseName}{Software-Defined Systems}
\newcommand{\StudentID}{6733172621}
\newcommand{\StudentName}{Patthadon Phengpinij}
\newcommand{\DefaultCollaborators}{Claude (for\,\LaTeX\,styling and grammar checking)}
\newcommand{\AssignmentLabel}{Activity}

\usepackage{../vendor/template/CEDT-Assignment-style}

% listings ships no HCL/Terraform dialect and the shared style only predefines
% YAML, so declare a minimal one here to colour main.tf the same way the YAML in
% the earlier activities is coloured.
\lstdefinelanguage{HCL}{
	morekeywords={terraform, provider, resource, variable, output, module, data,
		locals, required_providers, required_version, depends_on, true, false, null},
	sensitive=true,
	morecomment=[l]{\#},
	morecomment=[l]{//},
	morecomment=[s]{/*}{*/},
	morestring=[b]",
}

\begin{document}
\subject
\asgntitle{5}{}{Week 6}{\StudentID\ \StudentName}{\DefaultCollaborators}

% ================================================================================ %
\section*{Activity: DevOps and Infrastructure as Code}
% ================================================================================ %

\textbf{Instructions}
\begin{enumerate}
	\item This assignment is due Wednesday September 16, 2026 no later than 11:00 pm on mycourseville (pdf only).
	\item Use {terraform} to run \texttt{todo} container version 3.0 and \texttt{redis} container.
	\item Prepare a PDF report containing the following information:
	      \begin{enumerate}
		      \item Content of your \texttt{terraform} file.
		      \item Screenshot of running \texttt{terraform} in each step.
	      \end{enumerate}
\end{enumerate}

\begin{center}
	\par\noindent\textbf{\color{red} YOU MUST ALSO USING GRADER (URL is in MCV's assignment)}
\end{center}

\newpage

\noindent\textbf{Environment.}
Everything runs on a local Docker Engine 29.7.2 (Ubuntu on WSL2), driven by
\textbf{Terraform v1.13.3} through the \textbf{\texttt{kreuzwerker/\allowbreak docker}}
provider v3.9.0 --- the community provider that talks to the Docker Engine API
directly, so no \texttt{docker run} or \texttt{docker compose} is involved: the
containers are created, tracked in \texttt{terraform.tfstate}, and torn down
entirely as Terraform-managed resources.

\medskip

\noindent\textbf{A note on the third container.}
Version~3.0 of \texttt{todo} (\texttt{natawut/todo-service:release-3}) is the
release that gained notifications: every \texttt{POST /} calls
\texttt{set\_notification()}, which fires a request at
\texttt{NOTIFICATION\_HOST:\allowbreak NOTIFICATION\_PORT} (default \texttt{localhost:9000})
and never looks at the reply. The reply is ignored, but a \emph{refused}
connection is not: \texttt{requests} raises \texttt{ConnectionError} out of the
request handler, so with \texttt{todo} and \texttt{redis} alone every create
answers \texttt{500 Internal Server Error}, and the grader's POST check fails.
This activity has no notification service to point it at, so a plain
\texttt{nginx} stands in as that endpoint: it answers \texttt{405} to the
notification POST, \texttt{todo} discards the reply, and the create succeeds.
It is one \texttt{docker\_container} resource and one variable
(\texttt{var.notification\_image}); pointing that variable at the real
notification service is the only change needed to make it real.

\bigskip

% ================================================================================ %
%                                    Problem 01                                    %
% ================================================================================ %
\begin{problem}
Content of \texttt{terraform} file.
\end{problem}

\begin{solution}
	The whole configuration lives in one file, \texttt{main.tf}, and declares seven
	resources: the \texttt{todo-net} bridge network, three images, and three
	containers. Only \texttt{todo} is published to the host
	(\texttt{8000:8000}); \texttt{redis} and the notification stub stay inside
	\texttt{todo-net}, where Docker's embedded DNS resolves them by container name ---
	which is what \texttt{REDIS\_HOST=redis} relies on. The container is named
	\texttt{todo-service}, the name the rest of the course's stack (and the grader)
	expects. Declaring the images as \texttt{docker\_image} resources makes the pull
	part of the plan instead of a side effect of the run, and
	\texttt{keep\_locally~=~true} leaves the pulled images in the local cache when the
	rest is destroyed.
	\begin{codingbox*}[language=HCL]
# Activity 5 - Infrastructure as Code
# Run the todo service (version 3.0) and its redis database with Terraform.

terraform {
  required_version = ">= 1.5"

  required_providers {
    docker = {
      source  = "kreuzwerker/docker"
      version = "~> 3.0"
    }
  }
}

# Talks to the local Docker daemon. The host is left at the provider default so
# it follows DOCKER_HOST / unix:///var/run/docker.sock, which keeps the file
# portable between Docker Desktop and a plain dockerd.
provider "docker" {}

variable "todo_image" {
  description = "Docker image of the todo service (version 3.0)."
  type        = string
  default     = "natawut/todo-service:release-3"
}

variable "redis_image" {
  description = "Docker image of the redis database."
  type        = string
  default     = "redis:7-alpine"
}

# todo 3.0 POSTs a notification to NOTIFICATION_HOST:NOTIFICATION_PORT on every
# created todo and never checks the response, but a refused connection raises
# out of the request handler and turns POST / into a 500. This activity has no
# notification service to point it at, so a plain nginx stands in as the
# endpoint: it answers 405 to the notification POST, todo discards the response,
# and creating a todo succeeds. Swap this for the real notification service and
# nothing else in this file has to change.
variable "notification_image" {
  description = "Docker image standing in for the notification endpoint."
  type        = string
  default     = "nginx:alpine"
}

variable "todo_port" {
  description = "Host port the todo service is published on."
  type        = number
  default     = 8000
}

# ---------------------------------------------------------------------------
# Network
# ---------------------------------------------------------------------------

# A user-defined bridge network gives both containers Docker's embedded DNS, so
# todo can reach redis by container name (REDIS_HOST=redis below).
resource "docker_network" "todo_net" {
  name   = "todo-net"
  driver = "bridge"
}

# ---------------------------------------------------------------------------
# Images
# ---------------------------------------------------------------------------

resource "docker_image" "redis" {
  name         = var.redis_image
  keep_locally = true
}

resource "docker_image" "todo" {
  name         = var.todo_image
  keep_locally = true
}

resource "docker_image" "notification" {
  name         = var.notification_image
  keep_locally = true
}

# ---------------------------------------------------------------------------
# Containers
# ---------------------------------------------------------------------------

# Not published to the host: only the todo service needs to reach it, and it
# does so over todo-net.
resource "docker_container" "redis" {
  name    = "redis"
  image   = docker_image.redis.image_id
  restart = "unless-stopped"

  networks_advanced {
    name = docker_network.todo_net.name
  }
}

# Answers the notification POST described by var.notification_image. Like redis,
# it is internal to todo-net only.
resource "docker_container" "notification" {
  name    = "notification-service"
  image   = docker_image.notification.image_id
  restart = "unless-stopped"

  networks_advanced {
    name = docker_network.todo_net.name
  }
}

resource "docker_container" "todo" {
  name    = "todo-service"
  image   = docker_image.todo.image_id
  restart = "unless-stopped"

  env = [
    "REDIS_HOST=${docker_container.redis.name}",
    "REDIS_PORT=6379",
    "NOTIFICATION_HOST=${docker_container.notification.name}",
    "NOTIFICATION_PORT=80",
  ]

  ports {
    internal = 8000
    external = var.todo_port
  }

  networks_advanced {
    name = docker_network.todo_net.name
  }

  # Terraform already infers these from the REDIS_HOST / NOTIFICATION_HOST
  # references above; stated explicitly so the ordering survives any edit to
  # those lines.
  depends_on = [
    docker_container.redis,
    docker_container.notification,
  ]
}

# ---------------------------------------------------------------------------
# Outputs
# ---------------------------------------------------------------------------

output "todo_url" {
  description = "Where the todo service is reachable from the host."
  value       = "http://localhost:${var.todo_port}"
}

output "container_names" {
  description = "Containers created by this configuration."
  value = [
    docker_container.todo.name,
    docker_container.redis.name,
    docker_container.notification.name,
  ]
}
	\end{codingbox*}
\end{solution}
% ================================================================================ %

% ================================================================================ %
%                                    Problem 02                                    %
% ================================================================================ %
\begin{problem}
Screenshot of running \texttt{terraform} in each step.
\end{problem}

\begin{solution}
	\textbf{Step 1 --- \texttt{terraform init}.}
	Downloads the \texttt{kreuzwerker/docker} plugin into \texttt{.terraform/} and
	records the exact version in \texttt{.terraform.lock.hcl}.
	\begin{codingbox*}[language=bash]
$ terraform init

Initializing the backend...
Initializing provider plugins...
- Finding kreuzwerker/docker versions matching "~> 3.0"...
- Installing kreuzwerker/docker v3.9.0...
- Installed kreuzwerker/docker v3.9.0 (self-signed, key ID 0DCE698927DAF8EC)
Partner and community providers are signed by their developers.
[...]
Terraform has created a lock file .terraform.lock.hcl to record the provider
selections it made above. Include this file in your version control repository
so that Terraform can guarantee to make the same selections by default when
you run "terraform init" in the future.

Terraform has been successfully initialized!
	\end{codingbox*}

	\textbf{Step 2 --- \texttt{terraform plan}.}
	Seven resources to add, nothing to change or destroy. Only the
	\texttt{docker\_container.todo} block is reproduced here --- it is the one
	carrying the environment and the published port; the other six blocks and the
	attributes that are only known after apply are elided.
	\begin{codingbox*}[language=bash]
$ terraform plan

Terraform used the selected providers to generate the following execution
plan. Resource actions are indicated with the following symbols:
  + create

Terraform will perform the following actions:

  # docker_container.todo will be created
  + resource "docker_container" "todo" {
      + env                                         = [
          + "NOTIFICATION_HOST=notification-service",
          + "NOTIFICATION_PORT=80",
          + "REDIS_HOST=redis",
          + "REDIS_PORT=6379",
        ]
      + name                                        = "todo-service"
      + network_mode                                = "bridge"
      + restart                                     = "unless-stopped"
      [... remaining attributes known after apply ...]

      + networks_advanced {
          + aliases      = []
          + name         = "todo-net"
        }

      + ports {
          + external = 8000
          + internal = 8000
          + ip       = "0.0.0.0"
          + protocol = "tcp"
        }
    }

  [... docker_container.redis, docker_container.notification, docker_image.todo,
       docker_image.redis, docker_image.notification, docker_network.todo_net ...]

Plan: 7 to add, 0 to change, 0 to destroy.

Changes to Outputs:
  + container_names = [
      + "todo-service",
      + "redis",
      + "notification-service",
    ]
  + todo_url        = "http://localhost:8000"
	\end{codingbox*}

	\textbf{Step 3 --- \texttt{terraform apply}.}
	Terraform reprints the plan and then creates the resources, respecting the
	dependency order it worked out: the images and the network first, then
	\texttt{redis} and the notification stub, then \texttt{todo-service}, which
	depends on both. The pulls finish in a second because the images were already in
	the local cache from an earlier run --- \texttt{keep\_locally~=~true} is what
	keeps them there across a \texttt{destroy}.
	\begin{codingbox*}[language=bash]
$ terraform apply -auto-approve

[... same plan as above ...]

docker_image.notification: Creating...
docker_network.todo_net: Creating...
docker_image.redis: Creating...
docker_image.todo: Creating...
docker_image.todo: Creation complete after 1s [id=sha256:8564bc13a323...natawut/todo-service:release-3]
docker_image.redis: Creation complete after 1s [id=sha256:ff02b58f971e...redis:7-alpine]
docker_image.notification: Creation complete after 1s [id=sha256:72ba65eb42c1...nginx:alpine]
docker_network.todo_net: Creation complete after 2s [id=450a852a4874...]
docker_container.redis: Creating...
docker_container.notification: Creating...
docker_container.notification: Creation complete after 0s [id=fb81a2187631...]
docker_container.redis: Creation complete after 0s [id=3b64528cd7d9...]
docker_container.todo: Creating...
docker_container.todo: Creation complete after 1s [id=1e914733dee6...]

Apply complete! Resources: 7 added, 0 changed, 0 destroyed.

Outputs:

container_names = [
  "todo-service",
  "redis",
  "notification-service",
]
todo_url = "http://localhost:8000"
	\end{codingbox*}

	\textbf{Step 4 --- verifying what Terraform built.}
	All three containers are up, with port \texttt{8000} the only one published to
	the host, and the service does a full round trip: the \texttt{POST} returns the
	id of the new todo and the following \texttt{GET} reads it back out of redis.
	\begin{codingbox*}[language=bash]
$ docker ps --filter name=todo-service --filter name=redis --filter name=notification-service \
      --format "table {{.Names}}\t{{.Image}}\t{{.Status}}\t{{.Ports}}"
NAMES                  IMAGE          STATUS         PORTS
todo-service           8564bc13a323   Up 8 seconds   0.0.0.0:8000->8000/tcp
notification-service   72ba65eb42c1   Up 9 seconds   80/tcp
redis                  ff02b58f971e   Up 9 seconds   6379/tcp

$ curl -X POST -d '{"title":"todo-1","detail":"the first todo","duedate":"2022-10-24 11:00:00","tags":[],"completed":false}' http://localhost:8000
{"id":0}

$ curl http://localhost:8000
[{"id":0,"title":"todo-1","detail":"the first todo","completed":false,"duedate":"2022-10-24T11:00:00","tags":[]}]
	\end{codingbox*}
\end{solution}
% ================================================================================ %

% ================================================================================ %
%                                    Problem 03                                    %
% ================================================================================ %
\begin{problem}
Screenshot of the results after running the \texttt{GRADER}.
\end{problem}

\begin{solution}
	The grader is pointed at \texttt{main.tf} (its default), checks the Terraform
	side --- file present, CLI installed, working directory initialized, plan
	generated --- and then the running stack: both named containers, the service
	answering on \texttt{http://localhost:8000}, and a \texttt{POST} whose result
	comes back on the following \texttt{GET}. Every check passes ---
	\texttt{Result: true} --- after which it prompts for the Student ID and full name
	to submit the result. (In the transcript below the per-line timestamps are
	trimmed and the status glyphs the grader prints are written as \texttt{OK}.)
	
	\fig[\linewidth]{images/grader.png}{\texttt{activity5-CEDT} grader --- all checks pass, \texttt{Result: true}.}{grader}
\end{solution}
% ================================================================================ %

\end{document}
